\documentclass[11pt,a4paper]{article}
\usepackage[italian]{babel}
\usepackage[margin=2.5cm]{geometry}
\usepackage{parskip}
\usepackage{amsmath, amssymb, amsthm}
\usepackage[most]{tcolorbox}
\usepackage{array}
\usepackage{hyperref}

% --- Riquadri con bordo nero, senza numero ---
\tcbset{nota/.style={colback=white, colframe=black, boxrule=0.5pt,
  sharp corners}}
\NewTColorBox{definizione}{}{nota,
  before upper={\textbf{Definizione.}\ }}
\NewTColorBox{teorema}{o}{nota,
  before upper={\textbf{Teorema\IfValueT{#1}{ (#1)}.}\ }}
\NewTColorBox{proposizione}{o}{nota,
  before upper={\textbf{Proposizione\IfValueT{#1}{ (#1)}.}\ }}

% --- Ambienti semplici (senza riquadro) ---
\theoremstyle{definition}
\newtheorem*{esempio}{Esempio}
\newtheorem*{osservazione}{Osservazione}

% V e F nelle tavole di verità
\newcommand{\V}{\textsf{V}}
\newcommand{\F}{\textsf{F}}

\title{Logica}
\author{Marco Cerbella}
\date{Lezione 1 -- 21 settembre 2026}

\begin{document}
\maketitle

\section{Enunciati e connettivi}

\begin{definizione}
Un \emph{enunciato} (o \emph{proposizione}) è una frase di cui possiamo
dire con certezza se è \V\ (vera) oppure \F\ (falsa). Gli enunciati si
indicano con $p, q, r, \dots$
\end{definizione}

\begin{esempio}
Sono enunciati: ``$5 > 3$'', ``$99$ è divisibile per $3$'',
``$508$ è multiplo di $5$''. Invece ``$n$ è un numero primo'' \textbf{non}
è un enunciato, perché la sua verità dipende da $n$.
\end{esempio}

Dati due enunciati $p, q$ se ne costruiscono di nuovi con i
\emph{connettivi}:
\begin{center}
\begin{tabular}{cll}
$p \land q$ & ``$p$ e $q$'' & congiunzione \\
$p \lor q$ & ``$p$ o $q$'' & disgiunzione \\
$\neg p$ & ``non $p$'' & negazione \\
$p \to q$ & ``se $p$ allora $q$'' & implicazione \\
$p \leftrightarrow q$ & ``$p$ se e soltanto se $q$'' & coimplicazione
\end{tabular}
\end{center}

\subsection*{Tavola di verità}

\begin{center}
\begin{tabular}{cc|c|c|c|c}
$p$ & $q$ & $p \land q$ & $p \lor q$ & $p \to q$ & $p \leftrightarrow q$ \\ \hline
\V & \V & \V & \V & \V & \V \\
\V & \F & \F & \V & \F & \F \\
\F & \V & \F & \V & \V & \F \\
\F & \F & \F & \F & \V & \V
\end{tabular}
\end{center}

\section{Proprietà dei connettivi}

\begin{proposizione}
Siano $p, q, r$ enunciati. Allora:
\begin{enumerate}
    \item $p \lor q \iff q \lor p$;
    \item $p \land q \iff q \land p$;
    \item $(p \land q) \lor r \iff (p \lor r) \land (q \lor r)$;
    \item $(p \lor q) \land r \iff (p \land r) \lor (q \land r)$;
    \item $\neg(p \land q) \iff (\neg p) \lor (\neg q)$;
    \item $\neg(p \lor q) \iff (\neg p) \land (\neg q)$;
    \item $p \to q \iff (\neg p) \lor q$;
    \item $p \to q \iff \neg q \to \neg p$;
    \item $p \leftrightarrow q \iff (p \to q) \land (q \to p)$.
\end{enumerate}
\end{proposizione}

Due enunciati composti sono equivalenti ($\iff$) quando hanno la stessa
tavola di verità. Le prime due proprietà sono la commutatività di $\lor$ e
$\land$; la 3 e la 4 sono le proprietà distributive; la 5 e la 6 sono le
leggi di De Morgan; la 8 è la \emph{contronominale}.

\begin{proof}[Verifica della proprietà 7]
Confrontiamo le tavole di verità di $p \to q$ e di $(\neg p) \lor q$:
\begin{center}
\begin{tabular}{cc|c|c|c}
$p$ & $q$ & $p \to q$ & $\neg p$ & $(\neg p) \lor q$ \\ \hline
\V & \V & \V & \F & \V \\
\V & \F & \F & \F & \F \\
\F & \V & \V & \V & \V \\
\F & \F & \V & \V & \V
\end{tabular}
\end{center}
Le colonne di $p \to q$ e di $(\neg p) \lor q$ sono uguali.
\end{proof}

\begin{proof}[Verifica della proprietà 4]
Servono otto righe, una per ogni combinazione dei valori di $p, q, r$:
\begin{center}
\begin{tabular}{ccc|c|c|c|c|c}
$p$ & $q$ & $r$ & $p \lor q$ & $(p \lor q) \land r$ & $p \land r$ &
$q \land r$ & $(p \land r) \lor (q \land r)$ \\ \hline
\V & \V & \V & \V & \V & \V & \V & \V \\
\V & \V & \F & \V & \F & \F & \F & \F \\
\V & \F & \V & \V & \V & \V & \F & \V \\
\V & \F & \F & \V & \F & \F & \F & \F \\
\F & \V & \V & \V & \V & \F & \V & \V \\
\F & \V & \F & \V & \F & \F & \F & \F \\
\F & \F & \V & \F & \F & \F & \F & \F \\
\F & \F & \F & \F & \F & \F & \F & \F
\end{tabular}
\end{center}
Le colonne di $(p \lor q) \land r$ e di $(p \land r) \lor (q \land r)$
coincidono. Le altre proprietà si verificano allo stesso modo.
\end{proof}

\section{Formule aperte e quantificatori}

\begin{definizione}
Una \emph{formula aperta} è un'affermazione che contiene un'incognita, un
parametro o una variabile.
\end{definizione}

\begin{esempio}
\[
p(x):\ x + 3 = 0, \qquad
q(n):\ \text{``$n$ è un numero pari''}, \qquad
r(k):\ \text{``$k$ è multiplo di $5$''}.
\]
\end{esempio}

Una formula aperta non è né vera né falsa. Per passare dalla formula aperta
a un enunciato ci sono tre modi.

\begin{enumerate}
    \item \textbf{Sostituzione}: si sostituisce alla variabile un valore
          preciso. Ad esempio $p(3)$ è l'enunciato $3 + 3 = 0$ (falso),
          $p(-3)$ è $-3 + 3 = 0$ (vero), $q(8)$ è ``$8$ è un numero pari''
          (vero).

    \item \textbf{Quantificatore esistenziale} $\exists$ (``esiste''):
          \[
          \exists\, x \in \mathbb{R} : p(x)
          \]
          si legge ``esiste $x$ reale tale che $x + 3 = 0$'' ed è vero
          (basta $x = -3$). Invece
          $\exists\, x \in \mathbb{R} : x^2 = -10$ è falso.

    \item \textbf{Quantificatore universale} $\forall$ (``per ogni''):
          $\forall\, x \in \mathbb{R} : x + 3 = 0$ è falso, mentre
          $\forall\, x \in \mathbb{R} : x^2 \geq 0$ è vero.
\end{enumerate}

\begin{proposizione}[Negazione dei quantificatori]
\[
\neg\big(\forall x : p(x)\big) \iff \exists\, x : \neg p(x), \qquad
\neg\big(\exists\, x : p(x)\big) \iff \forall x : \neg p(x).
\]
Il secondo enunciato si scrive anche $\nexists\, x : p(x)$.
\end{proposizione}

Per negare una frase con un quantificatore si scambia $\forall$ con
$\exists$ e si nega la formula.

\section{Implicazione: condizioni necessarie e sufficienti}

L'implicazione $p \to q$ si legge in tre modi:
\begin{itemize}
    \item ``se $p$ allora $q$'';
    \item ``$p$ è condizione \emph{sufficiente} per $q$'';
    \item ``$q$ è condizione \emph{necessaria} per $p$''.
\end{itemize}

\begin{esempio}
Per ogni $n \in \mathbb{N}$: se $n$ è divisibile per $4$ (è $p$), allora
$n$ è pari (è $q$). Essere divisibile per $4$ è sufficiente per essere
pari; essere pari è necessario per essere divisibile per $4$.
\end{esempio}

\end{document}
